\documentclass[
  12pt,
  oneside,
  openright
]{abntex2}

% Idioma e Fonte (XeLaTeX)
\usepackage[brazil]{babel}
\usepackage{fontspec}
\setmainfont{Times New Roman}

\usepackage[alf]{abntex2cite}
\usepackage[a4paper,top=3cm,bottom=2cm,left=3cm,right=2cm]{geometry}

% Identação
\usepackage{indentfirst}

% Tabelas geradas pelo Pandoc
\usepackage{longtable}
\usepackage{booktabs}
\usepackage{array}
\usepackage{calc} % Essencial para o cálculo de colunas do Pandoc

% Imagens e Links (O Pandoc injeta \href e \includegraphics do Markdown)
\usepackage{graphicx}
\usepackage{hyperref}
\usepackage{needspace}

\let\oldsection\section
\renewcommand{\section}[1]{%
  \Needspace{8\baselineskip}%
  \oldsection{#1}%
}

\let\oldsubsection\subsection
\renewcommand{\subsection}[1]{%
  \Needspace{6\baselineskip}%
  \oldsubsection{#1}%
}

\let\oldsubsubsection\subsubsection
\renewcommand{\subsubsection}[1]{%
  \Needspace{4\baselineskip}%
  \oldsubsubsection{#1}%
}

% Pandoc pode gerar tabelas longtable sem legenda com:
% {\def\LTcaptype{none} % do not increment counter}
% No abntex2+hyperref isso pode disparar: "No counter 'none' defined".
% Definimos um contador dummy para satisfazer \LTcaptype=none.
\newcounter{none}
\providecommand{\theHnone}{none.\arabic{none}}

% Macro OBRIGATÓRIA do Pandoc para listas
\providecommand{\tightlist}{%
  \setlength{\itemsep}{0pt}\setlength{\parskip}{0pt}}

% Cabeçalho
\pagestyle{plain}

% Metadados injetados pelo frontmatter YAML do Markdown
\title{$title$}
\author{$author$}
\date{$date$}

\begin{document}

% O abntex2 já aplica espaçamento 1.5 padrão nativamente.
% A flag \textual garante o comportamento da formatação ABNT para o corpo.
\textual

% CAPA
\begin{capa}
  \centering
  {\ABNTEXchapterfont\large $author$ \par}
  \vfill
  {\ABNTEXchapterfont\bfseries\Large $title$ \par}
  \vfill
  {\large $date$}
\end{capa}

% FOLHA DE ROSTO
\begin{folhaderosto}
  \centering
  {\ABNTEXchapterfont $author$ \par}
  \vfill
  {\ABNTEXchapterfont\bfseries\Large $title$ \par}
  \vfill
  Trabalho apresentado como requisito parcial.  
  \vfill
  {\large $date$}
\end{folhaderosto}

% SUMÁRIO
\pdfbookmark[0]{\contentsname}{toc}
\tableofcontents*
\cleardoublepage

% CONTEÚDO (Injetado do inter.md)
$body$

\end{document}